%% notes-tex/008-xue-ffa-epistasis/sections/results.tex
%% [[experiments.008-xue-ffa.perspective-epistasis-in-metabolic-engineering]]
%% https://github.com/Mjvolk3/torchcell/tree/main/notes-tex/008-xue-ffa-epistasis/sections/results.tex
%%
%% SOURCE: every number below is printed by one of
%%   experiments/008-xue-ffa/scripts/perspective_figure_panels.py       (panels a-j)
%%   experiments/008-xue-ffa/scripts/notable_interaction_selection.py   (consensus panels)
%%   experiments/008-xue-ffa/scripts/ffa_epistatic_path_panels.py       (path enumeration)
%%   experiments/010-kuzmin-tmi/scripts/ffa_panel_model_score.py        (growth-model scoring;
%%       results/ffa_panel_model_score.{csv,json}, figure ffa_panel_model_score.svg)
%% Re-run those four and the panels and the numbers move together. Nothing here is
%% transcribed by hand from a figure.
%%
%% FIGURES. Each panel is a true-size PDF converted by `make plots` from the SVG its plot
%% script wrote. The same SVGs are embedded in notes/assets/drawio/ffa-epistasis-fig*.drawio,
%% which is where the arrangement is workshopped; \panelrow here is the layout that lets the
%% document compile before that arrangement is settled. Replace a \panelrow block with a
%% single \tcfig on the exported PDF once its figure is final in draw.io.

\section{Results\secstatus{todo}}

\subsection*{A complete three-factor design over ten transcription factors}

The dataset is a combinatorial deletion series in a free-fatty-acid-producing
\org{S.\ cerevisiae} strain. Ten transcription factors (FKH1, GCN5, MED4, OPI1, RFX1, RGR1,
RPD3, SPT3, TFC7 and YAP6) were deleted in every combination of one, two and three genes,
giving 10 single, 45 double and 120 triple deletion strains alongside the undeleted base
strain. The base strain is itself an engineered chassis, an FFA-overproducing
\org{S.\ cerevisiae} carrying pox1$\Delta$ faa1$\Delta$ faa4$\Delta$, and the ten
transcription factor deletions are additional to those three. Every titer below is relative
to that chassis, so a triple deletion strain carries six deletions in total. Six quantities
were measured per strain: the titers of myristate (C14:0), palmitate
(C16:0), stearate (C18:0), palmitoleate (C16:1) and oleate (C18:1), and their sum, which we
refer to as total titer. All titers below are normalized to the base strain, so a value of
1 is no change and 2 is a doubling.

Completeness is what makes the design usable for the question at hand. Because every
sub-combination of every triple was built, each triple can be reached by six distinct
deletion orders, and the intermediate strain at each step of each order was measured rather
than modeled. The design therefore contains 720 fully measured construction routes.

Only one of the ten single deletions raises titer at all. Deleting GCN5 gives
$1.06\times$; the other nine lose ground, and deleting RPD3 costs nearly half the titer at
$0.52\times$ (Fig.~\ref{fig:inaccessible}a). Combinations behave very differently: 29 of the
45 doubles and 51 of the 120 triples exceed the base strain, and the best strain in the
design is a triple at $2.05\times$. The gap between what single deletions do and what
combinations do is the first sign that the effects are not composing.

\begin{figure*}[t]
\panelrow{%
  \panel{a}{figures/panel_interaction_distribution.pdf}\hfill
  \panel{b}{figures/panel_volcano_total_titer.pdf}\hfill
  \panel{c}{figures/panel_consensus_agreement.pdf}%
}
\caption{\textbf{Interaction is the rule on total fatty acid titer, and the third deletion
is where it bites.} \textbf{a}, Interaction scores against the multiplicative null for all
45 double (lilac) and 120 triple (amber) deletions, as outline histograms on shared bins so
the heavy overlap stays readable. The median moves from $+0.40$ at two deletions to $-0.72$
at three. \textbf{b}, Trigenic interaction score $\tau$ against
evidence for the 119 testable triples under the multiplicative model. The dashed line marks
the largest $P$ that Benjamini-Hochberg rejects at a 5\% false discovery rate over the 119
tests of this readout ($P = 0.036$); 86 triples clear it and 11 of those have $\tau > 0$.
\textbf{c}, How many of the four epistasis models call each triple significant, split by
the sign of $\tau$. Seventy-three triples are called by all four with the sign agreeing,
and none of the 73 is positive.}
\label{fig:landscape}
\end{figure*}

\subsection*{Interaction is the rule, and the third deletion is where it bites}

We scored each combination against a multiplicative null, in which the expected phenotype of
a combination is the product of its parts' phenotypes. For a double deletion the digenic
score is $\varepsilon_{ij} = f_{ij} - f_i f_j$. For a triple, the trigenic score removes the
digenic contributions before asking what remains:
%
\begin{equation}
\tau_{ijk} = f_{ijk} - f_{ij} f_k - f_{ik} f_j - f_{jk} f_i + 2 f_i f_j f_k .
\label{eq:tau}
\end{equation}
%
This is the trigenic score of the yeast genetic interaction
maps~\citep{kuzminSystematicAnalysisComplex2018,baryshnikovaQuantitativeAnalysisFitness2010},
applied to titer in place of fitness. Substituting the digenic scores into
Eq.~\ref{eq:tau} rewrites it as a residual against the pairs,
%
\begin{equation}
\tau_{ijk} = f_{ijk} - f_i f_j f_k
- \left( \varepsilon_{ij} f_k + \varepsilon_{ik} f_j + \varepsilon_{jk} f_i \right) ,
\label{eq:tau-eps}
\end{equation}
%
an identity rather than an approximation. Read this way, $\tau$ is what the third deletion
adds once the singles and the pairs have been given their due, and each pairwise deviation
enters weighted by the remaining single-deletion phenotype. Standard errors propagate
through Eq.~\ref{eq:tau} by the delta method, and each $t$ statistic is referred to the
Welch-Satterthwaite effective degrees of freedom of that linear combination, which has a
median of 4.28 across the 714 testable trigenic scores.

Deviations from the null are large and are mostly negative at three deletions. The median
digenic score is $+0.40$ and the median trigenic score is $-0.72$
(Fig.~\ref{fig:landscape}a), so adding a third deletion does not damp the interaction, it
reverses its typical direction. Of the 119 triples with a testable standard error, 86 depart
from the null at a 5\% false discovery rate computed over the 119 tests of this readout, and
11 of those 86 depart upward (Fig.~\ref{fig:landscape}b).

The pairs are worth stating on their own, because they are what
Eq.~\ref{eq:tau-eps} measures the triples against. On total titer 43 of the 45 doubles
have a positive digenic score, 25 clear the same false discovery rate and all 25 of those
are positive, and 93 of the 120 triples produce less titer than their own singles and pairs
predict (\suppnoteref{note:digenic}). The reversal at the third deletion is therefore a
change of direction rather than the continuation of a trend.

\begin{figure*}[t]
\panelrow{%
  \panel{a}{figures/panel_improving_by_order.pdf}\hfill
  \panel{b}{figures/panel_greedy_walk.pdf}\hfill
  \panel{c}{figures/panel_path_accessibility.pdf}%
}
\panelrow{\panel{d}{figures/ffa_epistatic_path_panels_top.pdf}}
\caption{\textbf{Improving combinations exist, and almost none of them is reachable one
deletion at a time.} \textbf{a}, Percentage of combinations exceeding the base strain at
one, two and three deletions (bars, counts above), with the best titer at each order (lilac
line, right axis). One of ten single deletions improves titer; the best triple reaches
$2.05\times$. \textbf{b}, A greedy campaign (brick) takes the best strictly improving single
deletion each round and stops when none improves. It reaches $1.80\times$ after two
deletions and halts. Gray lines are all six deletion orders of the highest-titer triple
(RFX1, RPD3, YAP6; star), each of which passes through an intermediate below the base
strain. \textbf{c}, The 120 triples classified by whether they beat the base strain and
whether any of their six orders improves at every step. \textbf{d}, Trajectories for the
six highest-titer triples, one colored line per deletion order, on a shared vertical scale.
Rung 0 is the pox1$\Delta$ faa1$\Delta$ faa4$\Delta$ chassis and the three rungs after it
add transcription factor deletions. Each panel gives that triple's final titer $f_{ijk}$ and
its trigenic score $\tau_{ijk}$, and the badge marks the sign and significance of $\tau$.
None of the six has a monotone route. Error bars are standard errors of the strain means.}
\label{fig:inaccessible}
\end{figure*}

\subsection*{Improving combinations are not reachable one deletion at a time}

An interaction score says a combination is surprising. It does not say whether a campaign
could have found the combination, and that is the question a strain engineer is actually
asking. We therefore scored each of the 720 measured routes for monotonicity: a route is
monotone when every deletion in it leaves the strain strictly more productive than the
strain before it, which is the condition under which a build-and-screen campaign that discards
non-improving rounds would follow the route to its end.

Two of the 720 routes are monotone. Relaxing the criterion to allow a step whose loss is
within one standard error raises the count to 11 routes across 9 triples, so the result is
not an artifact of measurement noise sitting just the wrong side of zero. Counting by
destination rather than by route: 51 of the 120 triples beat the base strain, and 2 of those
51 sit at the end of a monotone route (Fig.~\ref{fig:inaccessible}c). The other 49
improvements exist in the design and are invisible to a stepwise campaign.

The obstacle is a valley rather than a plateau. For each triple we took its shallowest
route, meaning the one whose worst intermediate is highest, and measured how far that worst
intermediate falls below the base strain. The median shallowest route dips 31\% below the
starting titer, and the deepest dips 48\%. A campaign built to discard anything worse than
its parent will discard the route at the first step.

Simulating the campaign directly gives the same answer in one number. Starting from the base
strain and taking the best strictly improving single deletion each round, with full
visibility of every candidate at every round, the campaign deletes GCN5 ($1.06\times$), then
TFC7 ($1.80\times$), and then stops, because all eight remaining third deletions lower the
titer (Fig.~\ref{fig:inaccessible}b). The best strain in the design is $2.05\times$, and
reaching it requires passing through an intermediate at $0.91\times$ on its most forgiving
order and $0.52\times$ on its least. The simulated campaign is more capable than a real one
in every respect that matters, and it still recovers 88\% of the achievable titer while
stopping two rounds short of the answer.

The greedy trap has a second form worth naming, because it is not the same failure. One of
the two monotone routes ends at FKH1, GCN5, MED4 at $1.73\times$. A campaign following it
would improve at every round and finish with a triple. The simulated campaign does not find
it, because at the second round TFC7 ($1.80\times$) beat FKH1 ($1.53\times$) and the better
local step led to a dead end. Being reachable is necessary for a stepwise campaign to find a
combination, and it is not sufficient.

\begin{figure*}[t]
\panelrow{%
  \panel{a}{figures/panel_model_agreement.pdf}\hfill
  \panel{b}{figures/panel_graph_enrichment.pdf}%
}
\panelrow{%
  \panel{c}{figures/panel_scale_scatter.pdf}\hfill
  \panel{d}{figures/panel_scale_slope.pdf}\hfill
  \panel{e}{figures/panel_readout_sign.pdf}%
}
\caption{\textbf{What counts as an interaction depends on the null, on the scale and on the
readout.} \textbf{a}, Triples called significant at a 5\% false discovery rate by each
combination of the four models. Ninety-seven triples are called by at least one model and 73
by all four. \textbf{b}, Fold enrichment of seven interaction graphs among the significant
trigenic triples, per model, over all six readouts pooled at $P < 0.05$. Asterisks mark
$P < 0.05$ by Fisher's exact test; all five significant results are depletions.
\textbf{c}, Trigenic score on the linear scale against the same quantity on the log scale.
The upper-left quadrant is empty, so no combination that is negative on one scale is
positive on the other. \textbf{d}, The 27 triples with a positive linear-scale score,
carried onto the log scale; 15 lose the sign. \textbf{e}, Consensus interactions per
readout, split by sign. Thirty-one positive consensus interactions exist across the five
species, and none of them survives on their sum.}
\label{fig:models}
\end{figure*}

\subsection*{What counts as an interaction depends on the null and on the scale}

Equation~\ref{eq:tau} is one choice among several, and a claim that rests on it should say
how much of the claim is the choice. We fit three alternatives on the same normalized strain
means: an additive null, in which expected phenotypes add rather than multiply; a
generalized linear model with a log link and a Gamma family; and ordinary least squares on
log titer differenced against the base strain.

The four models agree on which combinations interact far more than they agree on the sign.
Ninety-seven triples are called by at least one model at a 5\% false discovery rate within
the readout, 73 by all four, and every one of the 73 is negative (Fig.~\ref{fig:models}a).
The disagreement falls almost entirely on the positive calls, and it falls along the scale:
the multiplicative and additive models report 11 and 6 positive triples on total titer, and
the two log-scale models report none.

That split is not a difference of power, it is a difference of estimand. The log-scale
three-way interaction $g_{ijk} - g_{ij} - g_{ik} - g_{jk} + g_i + g_j + g_k$, with
$g = \log f$, agrees with the linear-scale score in magnitude ($r = 0.91$, Spearman
$\rho = 0.84$) but not near zero. Of the 27 triples with a positive linear-scale score, 15
are non-positive on the log scale, while no negative linear-scale score becomes positive
(Fig.~\ref{fig:models}c,d). Positive three-way epistasis on fatty acid titer is a
linear-scale statement, and reporting one without naming the scale is reporting the modeling
choice as a result.

The readout matters as much as the scale. Thirty-one triples clear all four models with the
sign agreeing on at least one of the five individual species, and 20 of those are on C14:0,
where every consensus interaction is positive. None survives on the summed total
(Fig.~\ref{fig:models}e). Summing the five species is a modeling choice too, and it is the
choice that removes the positive signal. The five species behind the sum are not variations
on one pattern: the number of triples called within a species runs from 3 to 80, C14:0
reverses sign between the digenic and the trigenic layer, and the summed total sits at
$r = 0.91$ from C16:0 and C18:1 but at $r = -0.28$ from C14:0
(\suppnoteref{note:species}). The four expectations themselves, and what each one costs to
adopt, are set out in \suppnoteref{note:models}.

\subsection*{No interaction network on hand predicts which combinations interact}

If interactions were predictable from prior knowledge, a campaign could avoid the problem by
choosing combinations it expects to be additive. We tested that directly by asking whether
the interacting triples are enriched for connectivity in existing interaction graphs. Seven
graphs cover these ten genes with at least one connected triple: physical interactions,
regulatory interactions, genetic interactions, TFLink, and the coexpression, experimental
and database channels of STRING 12.0. For each graph and each model we compared the fraction
of significant triples whose three genes form a connected subgraph to the same fraction
among the non-significant triples.

Nothing is enriched. Across 28 graph-by-model tests, 5 reach $P < 0.05$ and every one of the
5 is a depletion rather than an excess (Fig.~\ref{fig:models}b). The largest excess anywhere
is the STRING database channel under the multiplicative model at 1.31-fold, which does not
reach significance ($P = 0.13$). Drawing the significant interactions onto the pathway
itself makes the same point visually: the interacting triples span the transcription factor
set with no visible relation to which factors regulate which fatty acid genes
(Fig.~\ref{fig:network}).

\begin{figure*}[t]
\tcfig{figures/ffa-epistasis-fig4-network-overlay.pdf}{\textbf{Trigenic interactions on total fatty acid
titer, drawn upstream of the pathway they act through.} The ten deleted transcription
factors are arranged in a circle at left. Edges within the circle are the 86 trigenic
interactions significant at a 5\% false discovery rate within the total-titer readout under
the multiplicative model, 75 negative (brick) and 11 positive (blue), each drawn as the
three edges of its triangle; edge width is the number of interactions the pair takes part
in, from 0.40 pt at 1 to 2.15 pt at 8. Every one of the 45 pairs is in at least one
negative interaction, so a pair that also carries a positive one would hide it: the two
signs of a pair are therefore drawn as two parallel lines, offset perpendicular to the pair
by half their combined width, and neither covers the other. Dotted arrows are regulatory links from a factor to a
pathway gene in the SGD regulatory graph or TFLink; three of the ten factors have none.
The columns to the right are the 13 core fatty acid pathway genes, the 64 reactions they
catalyze, the intermediates those reactions consume or produce, and the 31 measured
species, one node per compartment (peroxisome, ER membrane, lipid particle) with a bracket
per species. LPA, 1-acyl-sn-glycerol 3-phosphate; PA, phosphatidate, written as the two
acyl chains.}{fig:network}
\end{figure*}

\subsection*{A model trained on growth interactions does not predict the titer interactions}

The other place a campaign could look for a predictor is a model trained on measured
combinations rather than on a graph. The torchcell equivariant cell graph transformer,
described in a companion manuscript in preparation
(\file{paper/nature-biotech/editing.pdf}), is trained on the trigenic growth interaction
screens of Kuzmin et al.~\citep{kuzminSystematicAnalysisComplex2018,kuzminExploringWholegenomeDuplicate2020},
where each triple deletion is scored by colony fitness in a wild-type background. What that
model is measured against, and how it behaves on the additive nulls, is worked through
in~\citep{volkWhat010Trigenic2026}. All ten
factors here are in its gene space, and seven of them were deleted somewhere in its
training data; GCN5, RPD3 and TFC7 were not. We scored each of the 120 triples as a
three-gene deletion and asked whether the predicted growth interaction tracks the measured
titer interaction under the multiplicative model (Fig.~\ref{fig:model-score}). The two
quantities differ in readout and in background, so a null here does not separate a model
that cannot rank trigenic interactions from a growth interaction that does not carry over to
titer. The question is only whether this predictor, as it stands, would help choose
combinations.

Neither magnitude nor rank transfers. Across the 120 triples the ensemble mean of three
checkpoints correlates with the measured $\tau$ at $r = 0.16$ ($P = 0.08$) and Spearman
$\rho = 0.13$ ($P = 0.17$), and the measured $\tau$ shows no trend across quartiles of the
prediction (Kruskal-Wallis $P = 0.39$). The three checkpoints agree with each other ($r$
0.62 to 0.71) far better than any of them agrees with titer ($r$ 0.08 to 0.20), so the miss
is a property of what the model learned rather than of one training run. Restricting to
the 86 triples called at a 5\% false discovery rate, to the 84 free of the contested TFC7
locus, or to the 35 whose three factors were all deleted in training moves $r$ between 0.12
and 0.16 without moving it out of the noise, and the top quartile by predicted $\tau$ holds
no more called interactions than the rest (22 of 30 against 64 of 90).

The one signal is the sign. Of the 20 triples the model predicts positive, 9 measure
positive, against 18 of the 100 it predicts non-positive: 45\% against 18\% (Fisher
$P = 0.016$). The same split holds in the 84 triples without TFC7 ($P = 0.005$) and in the
35 fully trained ones ($P = 0.04$). Every $P$ in this subsection is nominal among ten
related tests on one set of 120 triples, so the sign association is a lead to test on a new
panel, not a result. Read with the graph result, the training data a useful predictor needs
are combinations measured on the readout of interest: a model trained on growth
combinations carries over to titer at most the direction of the interaction.

\begin{figure*}[t]
\tcfig{figures/ffa_panel_model_score.pdf}{\textbf{A growth-trained transformer does not
predict the titer interactions.} Three checkpoints of the torchcell equivariant cell graph
transformer, trained on the Kuzmin 2018 and 2020 trigenic growth screens, scored each
transcription factor triple as a three-gene deletion in a wild-type background; the measured
quantity is the multiplicative titer interaction $\tau$ of the same triple on the
pox1$\Delta$ faa1$\Delta$ faa4$\Delta$ chassis, total titer. ``FFA titer FDR'' is
Benjamini-Hochberg over the 119 testable triples within total titer. \textbf{a}, Coverage:
120 scorable, 35 with all three factors seen as deletions in training (7 of 10 factors), 84
free of the contested TFC7 locus. \textbf{b}, Ensemble-mean predicted growth $\tau$ against
measured titer $\tau$, Pearson $r = 0.16$, $P = 0.08$; the key's classes number 86, 33 and
1. \textbf{c}, The same as ranks, Spearman $\rho = 0.13$, $P = 0.17$. \textbf{d}, Measured
$\tau$ by quartile bin of predicted $\tau$, bar at the bin mean, Kruskal-Wallis $P = 0.39$.
\textbf{e}, Pearson (filled) and Spearman (open) correlations: between checkpoints on the
same triples; each checkpoint, the ensemble mean and the worst of three against titer; the
ensemble on three subsets. \textbf{f}, Measured $\tau$ split by predicted sign (blue at or
below zero, purple above; filled where the titer interaction is FDR-called), 45\% against
18\% measured positive, Fisher $P = 0.016$. Every $P$ is nominal among ten related tests on
one set of 120 triples.}{fig:model-score}
\end{figure*}
